\section{Introduction}\label{sec:intro}

Indian farmers sell most of their produce through regulated wholesale markets (\emph{mandis}).
Each mandi publishes a daily minimum, maximum and modal (most frequent) price per commodity and variety through
AGMARKNET, and these reports are mirrored on the Open Government Data platform. The reports cover thousands of markets and hundreds of commodities.
They are the most detailed public price signal available to a farmer deciding whether to sell today, wait,
or take produce to another market.

Forecasting these prices has a large literature (Section~\ref{sec:related}). Most of it studies a single market, a few markets of one state, or monthly national indices,
with deep or hybrid models evaluated on a single proportional hold-out. Naive baselines are rarely reported, and prediction intervals are almost never reported
\cite{paul2022brinjal,manogna2025deep,paul2025deep,kumar2021onion}.
Forecast-evaluation research warns that such protocols can overstate gains: simple benchmarks are hard to beat
\cite{makridakis2020m4,makridakis2022m5}, scale-free errors are needed when series differ in level
\cite{hyndman2006another}, and leakage through splitting and preprocessing is a common pitfall
\cite{hewamalage2022pitfalls,cerqueira2020evaluating}.
Studies of market-information services in India also show that price information alone often does not raise the
prices farmers receive \cite{fafchamps2012sms,mitra2018asymmetric}, although some interventions have
had large effects \cite{jensen2007digital,goyal2010information}. A farmer-facing forecast should therefore be
honest about its uncertainty and should not be presented as a trading signal.

This paper reports the design and evaluation of AG-DATA, a distributed system built on Scala, Spark, Kafka and HDFS that
forecasts mandi-level modal prices for 36 commodities and serves the forecasts, with conformal intervals whose coverage we measure and with cautious decision cards,
to an installable six-language web application. Our contributions are evidence-based and deliberately modest:

\begin{enumerate}
  \item \textbf{A leakage-resistant benchmark at national scale.} The benchmark spans
    \num{11845} mandi--commodity--variety series and \num{21.1}M report-level rows.
    It uses chronological train/validation/test periods with target-date purging, five naive and seasonal baselines,
    and four horizons. Results are reported with MASE, per-horizon Diebold--Mariano tests and per-series win rates
    (Sections~\ref{sec:protocol}--\ref{sec:results}).
  \item \textbf{An empirical finding about the loss function.} Between 17\% and 34\% of target price changes
    are \emph{exactly} zero, because mandis often re-quote an unchanged modal price. This makes the choice of
    training loss decisive under absolute-error criteria. At the next-report horizon, squared-loss ridge regression and random forests
    do not differ significantly from persistence, whereas absolute-loss gradient boosting is 2.9\% better.
    A per-series exponential smoother is as good as any learned model at 1--3 days (Section~\ref{sec:results}).
  \item \textbf{Measured, not assumed, uncertainty.} Split-conformal intervals with Mondrian volatility bins are
    calibrated on a validation year. Their test-period coverage is reported by horizon, by volatility bin and before and after a
    change in the publishing system (Section~\ref{sec:intervals}).
  \item \textbf{Data-integrity findings that affect any study of this source.} These are a grade relabel in January 2025
    and a market rename and volume collapse after the November 2025 migration of AGMARKNET. Left untreated, they fragment
    series and silently shrink the test set (Section~\ref{sec:data}).
  \item \textbf{An end-to-end, reproducible system.} The system has shared batch/stream canonicalisation that is tested for parity,
    a Kafka/Structured Streaming path with stated delivery semantics, an HDFS layout, an API, and a
    low-bandwidth progressive web app in six languages whose decision rule is back-tested
    (Sections~\ref{sec:system}--\ref{sec:app}).
\end{enumerate}

We do not claim superiority over published models: no prior study uses the same data, horizons and splits,
and Section~\ref{sec:comparison} states why each published number is context only.
